% ds9-rust user manual
% Build: xelatex manual.tex   (or pdflatex; xelatex preferred for the °σ→ glyphs)
\documentclass[11pt,a4paper]{article}

\usepackage{geometry}
\geometry{margin=22mm,headheight=24pt}

\usepackage{fontspec}
% Use widely-available system fonts (fontconfig) so the manual builds on
% bare TeX Live installs without TeX Gyre / Latin Modern OTF bundles.
% Pin to the OTF files explicitly — fontconfig matches the t1 (Type 1)
% Nimbus first, but xdvipdfmx can't embed Type 1 PostScript fonts.
\setmainfont{NimbusRoman-Regular.otf}[
  Path           = /usr/share/fonts/urw-base35/,
  BoldFont       = NimbusRoman-Bold.otf,
  ItalicFont     = NimbusRoman-Italic.otf,
  BoldItalicFont = NimbusRoman-BoldItalic.otf,
]
\setsansfont{NimbusSans-Regular.otf}[
  Path           = /usr/share/fonts/urw-base35/,
  BoldFont       = NimbusSans-Bold.otf,
  ItalicFont     = NimbusSans-Italic.otf,
  BoldItalicFont = NimbusSans-BoldItalic.otf,
]
\setmonofont{DejaVu Sans Mono}[Scale=0.92]

\usepackage{xcolor}
\definecolor{accent}{HTML}{1F2D5C}      % deep navy
\definecolor{accentSoft}{HTML}{4663AD}  % softer navy for rules
\definecolor{warm}{HTML}{E07A5F}        % coral
\definecolor{teal}{HTML}{2A9D8F}        % teal
\definecolor{paper}{HTML}{FAF7F2}       % paper background
\definecolor{ink}{HTML}{1B1B1F}         % near-black text
\definecolor{rule}{HTML}{C9C2B6}        % hairline grey

\usepackage[colorlinks=true,
            linkcolor=accent,
            urlcolor=teal,
            citecolor=accent]{hyperref}
\hypersetup{
  pdftitle={ds9-rust manual},
  pdfauthor={ds9-rust contributors},
  pdfsubject={A Rust + slint port of SAOImage DS9},
  bookmarksnumbered=true,
}

\usepackage{titlesec}
\titleformat{\section}
  {\Large\sffamily\bfseries\color{accent}}{\thesection}{0.6em}{}
  [{\color{accentSoft}\titlerule[0.4pt]}]
\titleformat{\subsection}
  {\large\sffamily\bfseries\color{accent}}{\thesubsection}{0.5em}{}
\titleformat{\subsubsection}
  {\normalsize\sffamily\bfseries\color{accent}}{\thesubsubsection}{0.4em}{}

\usepackage{fancyhdr}
\pagestyle{fancy}
\fancyhf{}
\fancyhead[L]{\textsf{\color{accentSoft}ds9-rust manual}}
\fancyhead[R]{\textsf{\color{accentSoft}\nouppercase{\leftmark}}}
\fancyfoot[C]{\textsf{\color{accentSoft}\thepage}}
\renewcommand{\headrulewidth}{0.4pt}
\renewcommand{\headrule}{\hbox to\headwidth{\color{rule}\leaders\hrule height \headrulewidth\hfill}}

\usepackage[most]{tcolorbox}
\tcbset{
  enhanced,
  fonttitle=\sffamily\bfseries,
  boxrule=0.4pt,
  arc=2pt,
  left=8pt, right=8pt, top=4pt, bottom=4pt,
}
\newtcolorbox{noteb}{colback=accent!4!paper,colframe=accent!50!paper,title={Note}}
\newtcolorbox{tipb}{colback=teal!6!paper,colframe=teal!55!paper,title={Tip}}
\newtcolorbox{warnb}{colback=warm!8!paper,colframe=warm!60!paper,title={Heads up}}

\usepackage{listings}
\lstdefinelanguage{shell}{
  morekeywords={cargo,echo,nc,ds9,curl,ssh,find,xargs},
  morecomment=[l]{\#},
  morestring=[b]",
}
\lstdefinelanguage{ipc}{
  morekeywords={frame,scale,cmap,bin,zoom,region,file,save,value,sextractor,
                samp,xmlrpc,3d,hdu,movie,help,quit,version,xpaaccess,
                load,open,png,fits,tiff,eps,wcs,limits,smooth,gaussian,
                boxcar,median,grey,heat,red,green,blue,rainbow,sls,hsv,cool,a,b,bb,
                sky,projection,crop,centroid,radial,mosaic,tile,depth,slice,
                max,sum,mean,resolve,vizier,ned,next,previous,new,delete,clear,
                in,out,fit,reset,match,blink,composite,rgb},
  sensitive=true,
}
\lstset{
  basicstyle=\ttfamily\small,
  keywordstyle=\color{accent}\bfseries,
  commentstyle=\color{teal}\itshape,
  stringstyle=\color{warm},
  backgroundcolor=\color{paper!50!white},
  frame=single,
  framesep=4pt,
  rulecolor=\color{rule},
  showstringspaces=false,
  breaklines=true,
  columns=fullflexible,
  keepspaces=true,
  language=shell,
}

\usepackage{array}
\usepackage{booktabs}
\usepackage{enumitem}
\setlist{nosep,leftmargin=*}

\renewcommand{\arraystretch}{1.15}

% ---- title page ----
\title{}
\author{}
\date{}

\begin{document}
\pagecolor{paper}
\color{ink}

\begin{titlepage}
\centering
\vspace*{4cm}
{\sffamily\Huge\bfseries\color{accent} ds9-rust}\\[6pt]
{\color{accentSoft}\rule{0.45\textwidth}{0.6pt}}\\[14pt]
{\sffamily\Large User Manual --- v0.1}\\[4pt]
{\sffamily\large\color{accentSoft} A Rust + slint port of SAOImage DS9}\\[28pt]
\begin{tcolorbox}[width=0.7\textwidth,
                  colback=accent!5!paper,colframe=accent!40!paper,
                  halign=center, boxrule=0.4pt, arc=2pt]
\sffamily\small
A pragmatic, OGFinder-style FITS viewer for working astronomers ---\\
multi-frame, region-aware, WCS-correct, scriptable.
\end{tcolorbox}
\vfill
{\sffamily\small\color{accentSoft}
Smithsonian Astrophysical Observatory's SAOImage DS9, reimagined in Rust.\\
\href{https://github.com/kjhan0606/rusty_ds9}{github.com/kjhan0606/rusty\_ds9}\\[4pt]
This document was typeset with \LaTeX{} (xelatex).}
\end{titlepage}

\renewcommand{\contentsname}{\sffamily\Large\color{accent} Contents}
\tableofcontents
\thispagestyle{empty}
\clearpage

% ============================================================
\section{Quick start}

\begin{lstlisting}[language=shell]
# build (release recommended for large mosaics + smooth / bin)
cargo build --release -p ds9-app

# launch — open files via the File menu
./target/release/ds9

# preload an image (with optional region file + catalog)
./target/release/ds9 deep_field.fits regions.reg sources.cat
\end{lstlisting}

\begin{noteb}
The binary is named \texttt{ds9}, not \texttt{ds9-app}. The crate name and
binary name diverge on purpose --- it lets you alias \texttt{ds9} on your
\texttt{PATH} without colliding with the upstream Tcl/Tk DS9 if you keep
both installed.
\end{noteb}

\paragraph{System requirements.} Rust $\geq$ 1.85, a working OpenGL stack
(slint defaults to a winit + femtovg backend), and an X11 / Wayland / macOS
Cocoa display. Optional: \texttt{source-extractor} (or \texttt{sextractor},
or \texttt{sex}) on \texttt{PATH} for the catalog wrapper; \texttt{lpr} for
printing.

\subsection{Supported inputs}

\begin{tabular}{>{\sffamily}lp{0.65\textwidth}}
\toprule
Image cube  & \texttt{.fits}, \texttt{.fit}, \texttt{.fts},
              \texttt{.fz} (Rice tile-compressed), NAXIS=2 or NAXIS=3 \\
Region file & DS9 \texttt{.reg} --- circle / box / ellipse / annulus / point /
              line / polygon, in \texttt{image} or WCS coordinates
              (\texttt{fk5} / \texttt{icrs} / \texttt{galactic}) \\
Catalog     & TSV, whitespace, CSV, SExtractor ASCII\_HEAD,
              minimal VOTable (\texttt{<TR>/<TD>}) \\
Colormap    & 256-stop user LUT (text \texttt{R G B} per row, 0--1 or 0--255) \\
\bottomrule
\end{tabular}

% ============================================================
\section{Core concepts}

\paragraph{Frames.} Every loaded image is a \emph{frame}. A frame owns its
own stretch, limits, colormap, regions, catalog, view zoom \& pan,
smoothing, binning, orientation, and (for cubes) 3-D projection mode. The
info bar's \textsc{frm} cell shows \texttt{<active>/<total>}.

\paragraph{The render pipeline} runs in this order, every refresh:

\begin{center}
\sffamily
\colorbox{accent!10!paper}{\strut\ bin\ } $\rightarrow$
\colorbox{accent!10!paper}{\strut\ smooth\ } $\rightarrow$
\colorbox{accent!10!paper}{\strut\ stretch + limits\ } $\rightarrow$
\colorbox{accent!10!paper}{\strut\ colormap\ (LUT-aware)\ } $\rightarrow$
\colorbox{accent!10!paper}{\strut\ orientation\ }
\end{center}

\paragraph{Coordinates.} Every mouse handler converts display $(x, y)$ to
FITS pixel $(x_{\text{fits}}, y_{\text{fits}})$ via
\texttt{display\_to\_fits()} so flips, 180° rotations, and crops never
leak into region geometry. WCS readouts use a TAN gnomonic projection
parsed from \texttt{CRPIX/CRVAL/CD\_*}.

% ============================================================
\section{The window}

\begin{tabular}{@{}>{\sffamily}lp{0.66\textwidth}@{}}
\toprule
Menubar      & File / Edit / View / Frame / Bin / Zoom / Scale / Color /
               Region / Catalog / Analysis / Help. \\
Mode toggles & \texttt{pan} / \texttt{edit} / \texttt{region} /
               \texttt{crosshair} (next to the info bar). \\
Info bar     & cursor x/y, pixel value, WCS readout, frame count, mode. \\
Sidebar      & catalog table, marker list, optional panner / magnifier. \\
Canvas       & the image. Drag to pan; click to interact per mode. \\
Colorbar     & vertical colorbar with the active stretch + limits. \\
\bottomrule
\end{tabular}

\subsection{Modes}

\begin{description}[font=\sffamily\bfseries\color{accent}]
\item[pan]       drag = pan; click on a region selects it.
\item[edit]      drag = pan or move a region; click empty canvas to drop a
                 small circle marker.
\item[region]    click drops the active region template (default circle);
                 the property editor (Region $\triangleright$ Property)
                 lets you tune the just-placed region.
\item[crosshair] click pins a session crosshair in WCS (or pixel) space.
                 The crosshair follows the active frame: switch frames and
                 the same on-sky position is re-projected.
\end{description}

% ============================================================
\section{Menubar reference}

The menubar is the canonical surface --- every feature is reachable
without an IPC client.

\subsection{File}

\begin{tabular}{@{}>{\sffamily}lp{0.66\textwidth}@{}}
\toprule
Open\dots         & file picker; loads into a new frame. \\
Save Image\dots   & PNG of the current canvas (after stretch/cmap). \\
Save FITS\dots    & BITPIX=$-32$ float export with WCS preserved. \\
Save TIFF\dots    & in-house TIFF writer (no libtiff dep). \\
Save EPS\dots     & PostScript dump of the colour-mapped raster. \\
Print\dots        & renders to a temp PNG then \texttt{lpr}'s it. \\
SAMP Send Image   & broadcasts the current FITS path over SAMP (see \S\ref{sec:samp}). \\
SAMP Send VOTable & broadcasts a saved VOTable to other SAMP clients. \\
Quit              & exit. \\
\bottomrule
\end{tabular}

\subsection{Edit}

\begin{tabular}{@{}>{\sffamily}lp{0.66\textwidth}@{}}
\toprule
Undo / Redo      & not yet --- noop placeholders. \\
Crop             & sets the visible region of the active frame
                   (zoom + pan adjusted; original FITS untouched).
                   \emph{Reset Crop} restores fit. \\
Preferences\dots & toggles the Preferences panel
                   (cmap / stretch / limits / smooth / bin / blink-ms).
                   \texttt{Apply} pushes onto the active frame;
                   \texttt{Save} persists to
                   \texttt{\textasciitilde/.config/ds9-rust/prefs.toml}. \\
\bottomrule
\end{tabular}

\subsection{View}

Toggles the panner, magnifier, coordinate-grid overlay, crosshair, info
panel, mode buttons, and colorbar. All toggles persist for the session.

\subsection{Frame}

\begin{tabular}{@{}>{\sffamily}lp{0.66\textwidth}@{}}
\toprule
New Frame / Delete Frame & opens picker / removes active. \\
Next / Previous          & wraps. \\
Match\dots               & broadcast active frame's view (zoom / pan /
                            cmap / scale --- per the matching lock) to
                            siblings. \\
Blink                    & cycles every 500\,ms. \\
RGB Composite            & frames 1, 2, 3 $\to$ R, G, B channels of a new
                            frame. \\
HDU Movie                & auto-cycles HDUs of the active file every 800\,ms. \\
Mosaic WCS               & WCS-aware reprojection of all open frames into
                            one (\S\ref{sec:mosaic}). \\
Tile Frames              & toggles a grid view of all open frames; click a
                            tile to make it active. \\
3D Slice\dots            & for NAXIS=3 cubes: cycles through planes one
                            click at a time. \\
3D Max Intensity         & projects max along the cube axis. \\
3D Sum                   & sum projection. \\
3D Mean                  & mean projection (NaN-safe). \\
Rotate 180° / Flip H / Flip V / Reset Orientation & display-time
                            transforms; coords stay native. \\
Lock Zoom / Pan / Color / Scale & scope of \texttt{Match\dots} and
                            \texttt{Blink}. \\
\bottomrule
\end{tabular}

\subsection{Bin}

\begin{tabular}{@{}>{\sffamily}lp{0.66\textwidth}@{}}
\toprule
1 / 2 / 4 / 8 / 16 / 32        & block-bin factor (active frame). \\
Average / Sum / Sub-sample     & reduction mode for the bin operator. \\
\bottomrule
\end{tabular}

\subsection{Zoom}

In, Out, Fit (canvas), Reset (1:1), or pick a fixed
$1\times / 2\times / 4\times / 8\times / 16\times$.

\subsection{Scale}

Stretches: \texttt{linear}, \texttt{log}, \texttt{sqrt}, \texttt{squared},
\texttt{asinh}, \texttt{sinh}.\\
Limits: \texttt{zscale}, \texttt{minmax}.

\subsection{Color}

Built-in colormaps: \texttt{grey, red, green, blue, a, b, bb, heat, cool,
rainbow, sls, hsv}. \emph{Load Custom\dots} reads a 256-row LUT
(\texttt{R G B} text, 0--1 or 0--255); \emph{Clear Custom} drops back to
the built-in.

\subsection{Region}

New / Load\dots / Save\dots / Delete Selected / Delete All / Info /
Property\dots (opens the editor for the currently selected region ---
edit centre, radius, rotation, colour, label, then Apply).

\subsection{Catalog}

Load\dots, Clear, Run SExtractor\dots (\S\ref{sec:sex}), Online Query\dots
(\S\ref{sec:online}), Info.

\subsection{Analysis}

\begin{tabular}{@{}>{\sffamily}lp{0.66\textwidth}@{}}
\toprule
Pixel Table\dots   & 11$\times$11 raw values around the cursor. \\
Statistics\dots    & frame-wide $N$, min, max, mean, median, $\sigma$,
                     and finite-pixel count. \\
Histogram\dots     & 128-bin log-y histogram, computed against the
                     current stretch limits. \\
Contour Levels\dots & 5-level overlay (sign-change between zscale low/high). \\
Smooth (cycle)     & cycles Gaussian $\sigma \in \{0, 2, 4, 8\}$ px. \\
Smooth Off         & disables. \\
Smooth Kind\dots   & switches between Gaussian, Boxcar, Median kernels. \\
Centroid           & flux-weighted $(x, y)$ + half-light radius for the
                     selected region. \\
Radial Profile     & 1-D azimuthal profile around the selected region;
                     opens a floating plot. \\
Projection         & line-pixel profile for a \emph{line} region. \\
\bottomrule
\end{tabular}

% ============================================================
\section{Mouse \& keyboard}

\begin{tabular}{@{}>{\sffamily}lp{0.66\textwidth}@{}}
\toprule
Pan                 & drag (pan mode); arrow keys nudge by 1 canvas pixel. \\
Zoom                & menu, scroll-wheel, or IPC \texttt{zoom in}/\texttt{out}/N. \\
Hover readout       & x/y, value, WCS in the info bar. \\
Select region       & click on it (any mode). \\
Move region         & edit-mode press-and-drag. \\
Drop region         & edit-mode click on empty canvas (default circle, $r=6$). \\
Pick catalog source & click near the symbol or click the row in the table. \\
Crosshair           & crosshair-mode click (WCS-pinned if frame has WCS). \\
\bottomrule
\end{tabular}

% ============================================================
\section{Working with regions}

DS9 region files are read and written in their canonical form. Image-coord
regions look like:

\begin{lstlisting}
# Region file format: DS9 version 4.1
image
circle(512.0, 256.5, 18.0) # color=cyan text={hot pixel}
box(120, 80, 40, 25, 0)
ellipse(720, 400, 50, 30, 35) # color=yellow
annulus(960, 540, 8, 20)
point(160, 160) # point=cross
line(50, 50, 200, 200) # color=magenta
polygon(100, 100, 200, 110, 180, 200, 90, 180)
\end{lstlisting}

\paragraph{Properties carried.} \texttt{color=}, \texttt{width=},
\texttt{text=\{...\}}, dash, font. The Region $\triangleright$ Property
dialog edits all of these for the selected region.

% ============================================================
\section{WCS-aware regions}

When the active frame has a WCS, DS9-style WCS regions load correctly. The
following all parse:

\begin{lstlisting}
fk5
circle(12:34:56.7, +25:18:12, 30")
ellipse(189.123, 45.678, 1.5', 0.8', 30)
box(10h12m05s, -27d05'13", 2', 1', 0)
polygon(189.1, 45.7, 189.15, 45.72, 189.13, 45.78)
\end{lstlisting}

\paragraph{Sizes.} \texttt{"} = arcsec, \texttt{'} = arcmin, \texttt{d} =
deg. Decimal-degree numbers without a suffix are also accepted.

\paragraph{No WCS?} The parser falls back to \texttt{image} semantics and
the status bar warns you so you can spot the silent reinterpretation.

% ============================================================
\section{Source catalogs}

\paragraph{Auto-detect.} The loader sniffs the first $\sim$256 bytes:

\begin{itemize}
\item \texttt{<VOTABLE>} or \texttt{<TABLE>} \,$\to$\, VOTable parser.
\item Leading \texttt{\#} on first non-blank line \,$\to$\, SExtractor.
\item Tab in the first row \,$\to$\, TSV (header row).
\item Comma in the first row \,$\to$\, CSV (header row, quoted strings OK).
\item Otherwise \,$\to$\, whitespace-separated, columns synthesised
      (\texttt{col1, col2, \dots}).
\end{itemize}

\paragraph{X / Y mapping.} The catalog overlay needs image coords. The
loader looks up \texttt{X\_IMAGE/Y\_IMAGE} (SExtractor), falling back to
\texttt{XWIN\_IMAGE/YWIN\_IMAGE}, then bare \texttt{X/Y} or \texttt{x/y}.

\subsection{Run SExtractor}\label{sec:sex}

\textsf{Catalog $\triangleright$ Run SExtractor\dots} shells out to the
external SExtractor binary on the active frame's source FITS, parses
\texttt{ASCII\_HEAD} output, and overlays the result. Defaults:
\texttt{DETECT\_THRESH=1.5}, \texttt{DETECT\_MINAREA=5},
\texttt{BACK\_SIZE=64}. Override per-key via \texttt{SEXTRACTOR\_OPTS}:

\begin{lstlisting}[language=shell]
SEXTRACTOR_OPTS="-DETECT_THRESH 3.0 -DEBLEND_MINCONT 0.0001" \
  ./target/release/ds9 image.fits
\end{lstlisting}

The IPC verb \texttt{sextractor} runs the same code path.

\subsection{Online query}\label{sec:online}

\textsf{Catalog $\triangleright$ Online Query\dots} opens a small panel
with three buttons:

\begin{tabular}{@{}>{\sffamily}lp{0.66\textwidth}@{}}
\toprule
Resolve & SIMBAD/Sesame name $\to$ (RA, Dec). Recenters the active frame
          on the resolved position; drops a crosshair. \\
VizieR  & cone search around the cursor / crosshair (radius in arcmin).
          Result loads as a VOTable catalog. \\
NED     & NED \texttt{objsearch} cone search (returns a bar-separated
          ASCII table). \\
\bottomrule
\end{tabular}

All requests use HTTPS via \texttt{ureq + rustls} (no system OpenSSL).

% ============================================================
\section{Multi-frame, Match, Blink, RGB}

Every \textsf{File $\triangleright$ Open\dots} pushes a new frame.
\textsf{Frame $\triangleright$ Next / Previous} cycles them; per-frame
view state (zoom, pan, stretch, limits, cmap, regions, catalog, smooth,
bin) is saved on switch and restored when you come back.

\paragraph{Match.} Broadcasts the active frame's state to siblings. The
\textsf{Lock Zoom / Pan / Color / Scale} toggles control which channels
are pushed. \texttt{Lock Color} pushes both built-in cmaps and any custom
LUT.

\paragraph{Blink.} Cycles all frames every 500\,ms. Combine with
\texttt{Lock Pan} for true side-by-side comparison.

\paragraph{RGB Composite.} \textsf{Frame $\triangleright$ RGB Composite}
takes frames 1--3 (whichever are loaded), zscale-stretches each, packs
them into the R/G/B channels of a new frame, and pushes it onto the
stack. Useful for tri-band difference imaging.

\subsection{Mosaic WCS}\label{sec:mosaic}

\textsf{Frame $\triangleright$ Mosaic WCS} reprojects every open frame
that has a WCS into a single output frame. The output WCS is the first
WCS-bearing frame's, shifted to span the bounding box of all corners.
The reprojector uses output-driven nearest-neighbour:
$\text{world}(out) \to \text{pix}(in)$ for each output pixel. Overlap is
averaged. Frames without a WCS are skipped (with a warning).

\subsection{Tile Frames}

Toggles a grid view of all open frames (cols $= \lceil\sqrt{n}\rceil$).
Each tile is a per-frame RGBA composite scaled with aspect preservation.
Click any tile to switch to that frame and exit tile mode.

% ============================================================
\section{Filters: smoothing \& binning}

\paragraph{Smoothing.} Gaussian, Boxcar, or Median kernel; per-frame.
\textsf{Analysis $\triangleright$ Smooth (cycle)} steps Gaussian
$\sigma$ through $0 \to 2 \to 4 \to 8$ px. \textsf{Smooth Off} resets.
\textsf{Smooth Kind\dots} switches kernel.

\paragraph{Binning.} \textsf{Bin $\triangleright$ N} block-bins the
displayed image $N \times N$. \textsf{Bin $\triangleright$
Average / Sum / Sub-sample} picks the reduction. Visualisation only --- the
underlying FITS pixel grid and WCS are untouched.

% ============================================================
\section{3-D cube rendering}

When the active frame loaded a NAXIS=3 cube, the full plane stack is held
on the side; the 2-D display starts on plane 1. Use the \textsf{Frame}
menu's 3D entries to project:

\begin{tabular}{@{}>{\sffamily}lp{0.66\textwidth}@{}}
\toprule
3D Slice\dots      & each click advances to the next plane (with wrap).
                     IPC \texttt{3d slice N} sets a specific 1-based index. \\
3D Max Intensity   & per-pixel max along the cube axis (NaN-safe). \\
3D Sum             & per-pixel sum (skips NaN). \\
3D Mean            & per-pixel mean (skips NaN, divisor is finite count). \\
\bottomrule
\end{tabular}

The chosen mode is sticky on the frame --- smoothing, binning, stretch,
and cmap all run downstream of the projection.

% ============================================================
\section{Analysis tools}

\subsection{Statistics}
Frame-wide finite count, $\min / \max$, mean, median, $\sigma$.

\subsection{Histogram}
128-bin, log-$y$, computed against the current stretch limits.

\subsection{Contour overlay}
5 levels via sign-change detection between zscale low/high. Re-rendered
on every refresh; orientation-aware.

\subsection{Pixel table}
$11 \times 11$ raw pixel values centred on the cursor.

\subsection{Centroid}
Flux-weighted centroid plus a rough half-light radius for the selected
region (typically a circle / box / ellipse).

\subsection{Radial profile}
Azimuthally averaged 1-D profile, drawn in a floating plot window.

\subsection{Projection}
For a \emph{line} region, samples the pixel values along the line and
plots them. Useful for slit-style profile checks.

\subsection{Crop}
\textsf{Edit $\triangleright$ Crop} sets the active frame's visible region
to the current view. \textsf{Reset Crop} returns to a full-frame fit. The
underlying FITS image is never modified.

% ============================================================
\section{Saving, exporting, printing}

\begin{tabular}{@{}>{\sffamily}lp{0.66\textwidth}@{}}
\toprule
Save Image\dots & PNG of the current canvas (RGBA, alpha untouched). \\
Save FITS\dots  & BITPIX=$-32$, WCS keywords preserved, BSCALE=1 / BZERO=0. \\
Save TIFF\dots  & in-house writer; no libtiff. \\
Save EPS\dots   & PostScript dump of the colour-mapped raster. \\
Save JPEG\dots  & not built in (returns a status message). \\
Print\dots      & renders to a temp PNG, then pipes via \texttt{lpr}. \\
\bottomrule
\end{tabular}

% ============================================================
\section{IPC protocol}\label{sec:ipc}

ds9-rust ships a line-based, XPA-flavoured IPC protocol. Two transports
listen by default:

\begin{itemize}
\item Unix-domain socket at
      \texttt{\$XDG\_RUNTIME\_DIR/ds9-rust-\$USER.sock}
      (or \texttt{/tmp/\dots} on macOS).
\item TCP loopback on a dynamic port (e.g. \texttt{127.0.0.1:54861}).
\end{itemize}

The chosen paths/ports are written to
\texttt{\textasciitilde/.ds9-rust/xpa.info} so external clients can
discover them:

\begin{lstlisting}
pid=12345
unix=/run/user/1000/ds9-rust-alice.sock
tcp=54861
\end{lstlisting}

\subsection{Verb grammar}

\begin{tabular}{@{}>{\sffamily}lp{0.62\textwidth}@{}}
\toprule
quit, exit                & terminate. \\
xpaaccess, version        & XPA-style probes. \\
mode M                    & set window mode (pan / edit / region / crosshair). \\
frame N, frame next/prev  & switch frame. \\
frame new / delete        & menu equivalents. \\
frame mosaic / tile       & menu equivalents. \\
scale linear|log|sqrt|... & set stretch. \\
scale limits LO HI        & set user limits. \\
cmap NAME                 & switch built-in colormap. \\
bin N                     & block-bin factor. \\
zoom in|out|fit|reset|N   & view zoom. \\
pan to X Y                & pan in image coords. \\
region load / save PATH   & DS9 .reg I/O. \\
file open PATH            & load a FITS file (new frame). \\
save png|fits|tiff|eps P  & export. \\
value                     & emit the pixel value at the current cursor. \\
sextractor                & run the SExtractor wrapper. \\
samp image / votable PATH & broadcast over SAMP. \\
hdu next / list / N       & HDU navigator. \\
movie on / off            & toggle HDU movie. \\
3d max / sum / mean       & cube projection. \\
3d slice N                & cube slice index (1-based). \\
3d depth                  & query cube depth. \\
help                      & one-line summary. \\
\bottomrule
\end{tabular}

\subsection{Examples}

\begin{lstlisting}[language=shell]
SOCK=$(awk -F= '/^unix=/{print $2}' ~/.ds9-rust/xpa.info)

echo "frame next"             | nc -U "$SOCK"
echo "scale log"              | nc -U "$SOCK"
echo "cmap heat"              | nc -U "$SOCK"
echo "region load /tmp/x.reg" | nc -U "$SOCK"
echo "save png /tmp/out.png"  | nc -U "$SOCK"
echo "3d max"                 | nc -U "$SOCK"
echo "3d slice 12"            | nc -U "$SOCK"
echo "value"                  | nc -U "$SOCK"
echo "help"                   | nc -U "$SOCK"
\end{lstlisting}

The TCP transport accepts the same grammar:

\begin{lstlisting}[language=shell]
PORT=$(awk -F= '/^tcp=/{print $2}' ~/.ds9-rust/xpa.info)
echo "version" | nc 127.0.0.1 "$PORT"
\end{lstlisting}

\begin{warnb}
TCP listens only on \texttt{127.0.0.1}. Don't expose ds9-rust on a
public interface --- there's no auth on the IPC channel.
\end{warnb}

% ============================================================
\section{SAMP messaging}\label{sec:samp}

\textsf{File $\triangleright$ SAMP Send Image / VOTable} broadcasts the
current FITS path (or a saved VOTable) to other SAMP clients running on
the same desktop --- TOPCAT, Aladin, etc.

\paragraph{How it works.} ds9-rust reads the standard SAMP lockfile at
\texttt{\textasciitilde/.samp}, registers a client with the hub via
XML-RPC, then sends \texttt{samp.hub.notifyAll} with the canonical
\texttt{image.load.fits} or \texttt{table.load.votable} mtype. The XML-RPC
serialiser is hand-rolled (no SOAP / glib dep). On exit the client
unregisters cleanly.

\paragraph{Caveats.} Outbound only; ds9-rust does not currently subscribe
to incoming SAMP messages.

% ============================================================
\section{Preferences}

\textsf{Edit $\triangleright$ Preferences\dots} toggles a panel. Editable:
default colormap / stretch / limits / smoothing string / bin factor /
blink interval (ms). Buttons:

\begin{itemize}
\item \textbf{Apply} --- push onto the active frame.
\item \textbf{Save}  --- persist to \texttt{\textasciitilde/.config/ds9-rust/prefs.toml}
       (a simple \texttt{key = value} text file, not real TOML).
\item \textbf{Reset} --- restore built-in defaults.
\end{itemize}

The prefs file is loaded automatically at startup; new frames inherit
those defaults.

% ============================================================
\section{Cookbook}

\subsection*{Open a deep-field stack and bring up SExtractor sources}

\begin{lstlisting}[language=shell]
./target/release/ds9 deep_field.fits
# Catalog -> Run SExtractor...   (uses defaults; tweak SEXTRACTOR_OPTS if needed)
# Click a star on the image -> table jumps; click a row -> view recenters.
\end{lstlisting}

\subsection*{RGB difference image from three exposures}

\begin{lstlisting}[language=shell]
SOCK=$(awk -F= '/^unix=/{print $2}' ~/.ds9-rust/xpa.info)
echo "file open r.fits" | nc -U "$SOCK"
echo "file open g.fits" | nc -U "$SOCK"
echo "file open b.fits" | nc -U "$SOCK"
echo "frame 1; cmap red"   | nc -U "$SOCK"
echo "frame 2; cmap green" | nc -U "$SOCK"
echo "frame 3; cmap blue"  | nc -U "$SOCK"
# Frame -> RGB Composite  (creates a new packed frame)
\end{lstlisting}

\subsection*{Walk through a spectral cube}

\begin{lstlisting}[language=shell]
./target/release/ds9 cube.fits
# Frame -> 3D Slice...   (each click advances)
# Or via IPC:
echo "3d slice 17" | nc -U "$SOCK"
echo "3d max"      | nc -U "$SOCK"   # then back to a max-intensity view
\end{lstlisting}

\subsection*{Mosaic four CCD chips into one frame}

\begin{lstlisting}[language=shell]
./target/release/ds9 chip1.fits chip2.fits chip3.fits chip4.fits
# Frame -> Mosaic WCS
# (output frame uses chip1's WCS, shifted to fit all chip corners)
\end{lstlisting}

% ============================================================
\section{Troubleshooting}

\begin{tabular}{@{}p{0.4\textwidth}p{0.55\textwidth}@{}}
\toprule
\textsf{No matching fbConfigs / GLXBadFBConfig} on launch &
slint's OpenGL backend can't find a config. Try
\texttt{SLINT\_BACKEND=software ./target/release/ds9}, or update Mesa /
your GPU driver. \\
\textsf{RefCell already borrowed} when opening a file dialog (macOS) &
fixed in 4fa677f. If you're on an older build, pull \texttt{main}. \\
\textsf{file contains no 2-D image HDU} &
The file is a pure binary table or its primary HDU has \texttt{NAXIS=0}.
Use \textsf{Frame $\triangleright$ HDU Movie} or the HDU navigator
(sidebar) to step through extensions. \\
\textsf{tile-compressed .fz precision loss} &
Make sure the workspace \texttt{[patch.crates-io]} for vendored
\texttt{fitsrs} is intact. The patch upgrades unquantize to f64. \\
\textsf{SAMP Send Image silently does nothing} &
No SAMP hub running; start TOPCAT or Aladin first. \\
\bottomrule
\end{tabular}

% ============================================================
\section{Architecture}

\begin{tabular}{@{}>{\sffamily}lp{0.66\textwidth}@{}}
\toprule
ds9-app     & slint UI shell; all event wiring; menubar handler;
              frame / state management; IPC dispatch;
              mosaic + tile + 3D + SAMP + online-query glue. \\
ds9-fits    & FITS I/O wrapping vendored \texttt{fitsrs}; BITPIX
              8/16/32/$-32$/$-64$, BSCALE/BZERO/BLANK; tile-compressed
              \texttt{.fz}; minimal TAN-projection WCS; NAXIS=3 cube
              loader. \\
ds9-image   & stretches (linear/log/sqrt/squared/asinh/sinh), limits
              (zscale/minmax), built-in + custom colormaps, RGBA
              renderer, smoothing kernels, block-bin operators. \\
ds9-marker  & DS9 \texttt{.reg} parser/writer + region/marker model
              (image-coord shapes, with WCS round-trip handled by
              ds9-app). \\
ds9-catalog & TSV / SExtractor / CSV / VOTable parsers, column lookup,
              numeric sort. \\
vendor/fitsrs & local \texttt{fitsrs} patch (f64 unquantize for
              \texttt{.fz}). \\
\bottomrule
\end{tabular}

\paragraph{Persistent files} written by ds9-rust:

\begin{itemize}
\item \texttt{\textasciitilde/.config/ds9-rust/prefs.toml} --- preferences.
\item \texttt{\textasciitilde/.ds9-rust/xpa.info} --- IPC discovery
      (pid, unix socket path, tcp port).
\end{itemize}

\vfill
\begin{center}
\sffamily\small\color{accentSoft}
ds9-rust 0.1 --- MIT OR Apache-2.0 --- Smithsonian Astrophysical
Observatory's SAOImage DS9, in Rust.
\end{center}

\end{document}
